% !PRINTED 505
\weberpage{505}
%<<CONT
\gesperrt{in which the coefficients $A_1, A_2, \ldots, A_m$ are integral rational functionals.}

Functionals which do not belong among the integral ones we shall occasionally also, for the sake of brevity, designate as fractional functionals.

In justification of this definition we first prove the theorem:

3. \gesperrt{If $\omega$ is an integral functional according to Definition 2. and at the same time rational, then it is an integral rational functional (according to Definition 1.).}

Suppose that $\omega$ is a fractional rational functional, and hence that the absolute value of $\omega$ is a rational fraction $p:q$, in which $p$ and $q$ are integral rational numbers without a common divisor; then $q\omega$ is an integral rational functional whose absolute value is $=p$, and therefore relatively prime to $q$. But if, on the other hand, $\omega$ is at the same time integral in the sense of Definition 2., then we can set
\[
\omega^m = -(A_1\omega^{m-1} + A_2\omega^{m-2} \ldots)
\]
from which
\[
(q\omega)^m = -q[A_1(q\omega)^{m-1} + A_2\,q(q\omega)^{m-2} + \cdots) % [sic] opening square bracket is closed by a round parenthesis as printed
\]
follows.

But from this it results that the absolute value $p^m$ of $(q\omega)^m$ must be divisible by $q$, which is possible only if $q=1$. Definition 1. is therefore contained in Definition 2. as a special case.

If
\[
\Phi(t) = t^m + A_1t^{m-1} + A_2t^{m-2} + \cdots + A_m
\]
is a function of $t$ in which the coefficients $A_1, A_2, \ldots, A_m$ are fractional rational functionals with the variables $x, y, \ldots$, but free of the variable $t$, then one can determine a function $aE(x,y,\ldots)$ in which $a$ denotes the common denominator of the absolute values of $A_1, A_2, \ldots$ and $E(x,y,\ldots)$ a primitive integral function, so that
\weq{1}{aE(x\ y,\ldots)\,\Phi(t) = P(t,x,y,\ldots)} % [sic] printed without comma between x and y
is itself a primitive integral function [§.\,136, 2.].

If $\Phi(t)$ splits into two factors $\Phi_1(t)$, $\Phi_2(t)$ in $R$, that is, if
\[
\Phi(t) = \Phi_1(t)\,\Phi_2(t),
\]
then let one determine accordingly for the two functions $\Phi_1, \Phi_2$ the factors $a_1E_1, a_2E_2$, so that
%CONT>>
